% --- Table 1: SFT ---
\begin{table*}[h]
\centering
\scriptsize
\begin{tabular}{l|p{.89\textwidth}}
\textbf{RL} & \textbf{SFT Output} \\\hline
1 & . These have to brain with for to. The, the brain, and the. This and to. These and, and
the cells. The to a and to the heart that cells to is cells, are the their their. In, it
the heart of the brain this. The. Next of are cells the brain the heart the and by. But
the that the cells to to this the et al to the cells. These, and to the environment, a.
We the mechanism. The, and in,. In the body and are is. This cells, the body, and of the
heart to the body to of a a of, the brain. The, the. \\\hline
3 & brain a a of to the to, they. However. \\\hline
5 & , the that the cells to the of the the of of the visual in and of the brain are, the
the body. The. In, but the time to the eyes the brain and the brain cells. The their are
to the brain. Here, and the cells. \\\hline
7 & enerative their body to, the two the, in the body of, the their into. We from the cell
of, the into the by their cells. \\\hline
9 & brain. In. \\\hline
\end{tabular}
\caption[Outputs generated by SFT model at each requested Readability Level.]{Outputs generated by SFT model at each requested Readability Level (RL). Readability level is on a scale of $1$ to $9$, with $9$ being the most readable.}
\label{tab:sft-examples}
\end{table*}

% --- Table 3: DSPy ---
\begin{table}[h]
\centering
\scriptsize
\begin{tabular}{l|p{.89\textwidth}}
\textbf{RL} & \textbf{DSPy Output} \\\hline
1 & \textquotedbl\textquotedbl \\\hline
3 & \textquotedbl\textquotedbl \\\hline
5 & \textquotedbl\textquotedbl \\\hline
7 & \textquotedbl\textquotedbl \\\hline
9 & \textquotedbl\textquotedbl \\\hline
\end{tabular}
\caption[Outputs generated by DSPy at each requested Readability Level.]{Outputs generated by DSPy at each requested Readability Level (RL), empty because of length failure. Readability level is on a scale of $1$ to $9$, with $9$ being the most readable.}
\label{tab:dspy-examples}
\end{table}


% --- Table 2: LoRA ---
\begin{table*}[h]
\centering
\scriptsize

\begin{tabular}{l|p{.89\textwidth}}
\textbf{RL} & \textbf{LoRA Output} \\\hline
1 & Small open reading frames (smORFs) are small stretches of DNA that code for proteins,
but their translation into proteins has been difficult to detect. This is because they are
too short to produce the high conservation scores that are used to predict which stretches
of DNA code for proteins. To overcome this problem, Kondziella et al.\ have developed a
new method called Poly-Ribo-Seq that uses a combination of biochemistry and sequencing to
detect smORFs that are being translated into proteins. Using this method, they were able
to detect smORFs in \textit{Drosophila} cells that were not detected by other methods.
These smORFs were found to be translated at similar levels to canonical proteins and had
the same characteristics as canonical proteins. These findings suggest that many smORFs in
\textit{Drosophila} are being translated into proteins and may have important functions in
the cell. Future experiments will be needed to confirm this and to determine the functions
of these proteins. \\\hline
3 & The discoidin domain receptors (DDRs) are a family of receptor tyrosine kinases (RTKs)
that bind to collagen and play a key role in tissue development, cell migration, and
disease progression. DDRs are composed of an extracellular region that binds to collagen,
a transmembrane domain and a cytoplasmic region that contains the catalytic kinase domain.
Unlike other RTKs, DDRs are constitutively dimerised in the absence of ligand and activate
downstream signalling through a mechanism that is not well understood. Here, we show that
DDR1 activation by collagen involves phosphorylation of the juxtamembrane and activation
loop tyrosines in trans between dimers. Phosphorylation between DDR1 dimers requires
specific contacts between the transmembrane domains but not between the extracellular
domains. These findings define a unique activation mechanism for DDR1. \\\hline
5 & Stem cell biology is a rapidly evolving field that is critical for our understanding of
how tissues and organs form during development and how they are maintained in the adult.
In the fruit fly \textit{Drosophila melanogaster}, the lymph gland is a multi-lobed organ
that produces blood cells throughout the larval stages. This organ contains a stem cell
niche that maintains a pool of undifferentiated blood cell progenitors. While the adult
hematopoietic hub is known to be the final site of hematopoiesis, the mechanism of blood
cell formation in larvae has been largely unexplored. Here we identify a population of
cells in the early larval lymph gland that express the Notch signaling pathway. These
cells are multipotent, dividing and give rise to the domeless expressing blood cell
progenitors that populate the lymph gland. Notch signaling is required for their
maintenance and relies on Dpp signaling from the niche. These results reveal the presence
of hematopoietic stem cells in the early larval lymph gland and establish
\textit{Drosophila} as a model for studying hematopoietic stem cell biology. The findings
also shed light on the evolutionary conservation of hematopoietic stem cell biology
between \textit{Drosophila} and vertebrates. \\\hline
7 & The malaria parasite \textit{Plasmodium} causes a disease that kills more than a
million people every year. To infect the human body, the parasite must first invade the
liver cells. How this invasion takes place is still not well understood. Silvie et al.\
have now investigated how \textit{Plasmodium} parasites enter human liver cells. The
experiments show that two host proteins, CD81 and SR-BI, can be used by the parasite to
invade the liver cells. The parasite uses different proteins to enter the liver cells of
different species. The experiments also show that a parasite protein called P36 can bind
to the host proteins and trigger the invasion of the liver cells. The experiments provide
new insights into the invasion process and could ultimately help to develop new treatments
for malaria. \\\hline
9 & A large family of proteins called TRIM5 is thought to be a key part of the body's
natural defense against viruses. These proteins prevent certain viruses from multiplying
by binding to the proteins that form the outer shell of the virus. To find out how TRIM5
proteins work, researchers have made a detailed 3D image of a TRIM5 protein bound to a
virus. The image shows that the protein forms a hexagonal pattern around the virus, with
each hexagon made up of 12 smaller protein molecules. The next step will be to find out
how the proteins interact with the virus to prevent it from multiplying. These findings
could also help to develop new treatments for viruses. \\\hline
\end{tabular}
\caption[Outputs generated by LoRA model at each requested Readability Level.]{Outputs generated by LoRA model at each requested Readability Level (RL). Readability level is on a scale of $1$ to $9$, with $9$ being the most readable.}
\label{tab:lora-examples}
\end{table*}



\begin{table*}[]
\centering
\scriptsize
\begin{tabular}{l|p{.89\textwidth}}
\textbf{RL} & \textbf{ICL Output} \\ \hline
0 & The study aimed to investigate the early development of the zebra finch (Taeniopygia guttata) and its potential use as a model organism for studying embryogenesis and stem cell biology. The researchers characterized the morphology and molecular features of finch embryos at oviposition, which were found to be equivalent to the blastocyst stage mouse embryo. They also identified the expression of pluripotency-associated genes, including Nanog, PouV, and Dnmt3b, in the finch blastoderm at oviposition.  The study used quantitative PCR (Q-PCR) to compare the expression levels of various genes associated with early embryonic development and pluripotency in finch and chick blastoderms at oviposition. The results showed that finch blastoderms expressed higher levels of naive pluripotency markers, such as Fbxo15, Prdm14, and Tbx3, compared to chick blastoderms. \texttt{[...]}
% The finch blastoderm was also found to be more similar to the mouse naive state, which is characterized by the expression of these naive pluripotency markers.  The researchers also cultured finch and chick blastodermal cells in the presence of LIF and the MEK inhibitor, PD0325901, and found that finch cells formed aggregates that showed strong AP activity and expressed Nanog and PouV, while chick cells did not. This suggests that finch blastodermal cells may act more like mESCs, while chick blastodermal cells behave more similarly to mEpiSCs.  The study's findings have implications for our understanding of embryogenesis and stem cell biology in avians and mammals. The zebra finch may be a useful model organism for studying early embryonic development and stem cell biology, and the results of this study may contribute to the development of new technologies for understanding and manipulating avian development.  The study also highlights the importance of considering the evolutionary conservation of developmental processes across species. The findings suggest that the naive state of embryogenesis is conserved amongst amniotes, and that the pre-epithelialized epiblast in amniotes may be a key feature of this stage.  Overall, the study provides new insights into the early development of the zebra finch and its potential use as a model organism for studying embryogenesis and stem cell biology. 
The findings have implications for our understanding of developmental biology and may contribute to the development of new technologies for understanding and manipulating avian development. \\ \hline
25 & The study aimed to investigate the early development of the zebra finch (Taeniopygia guttata) and its potential use as a model organism for studying embryogenesis and stem cell biology. The researchers characterized the morphology and molecular features of finch embryos at oviposition, which were found to be equivalent to the blastocyst stage mouse embryo. They also identified the expression of pluripotency-associated genes, including Nanog, PouV, and Dnmt3b, in the finch blastoderm at oviposition.  \texttt{[...]}
% The study used quantitative PCR (Q-PCR) to compare the expression levels of various genes associated with early embryonic development and pluripotency in finch and chick embryos.
% The results showed that finch blastoderms at oviposition expressed higher levels of naive pluripotency markers, such as Fbxo15, Prdm14, and Tbx3, compared to chick blastoderms. 
% The finch blastoderm was also found to be more similar to the mouse naive state, which is characterized by the expression of these naive pluripotency markers.  The researchers also cultured finch and chick blastodermal cells in the presence of LIF and the MEK inhibitor, PD0325901, and found that finch cells formed aggregates that showed strong AP activity and expressed Nanog and PouV, while chick cells did not. This suggests that the finch blastoderm at oviposition is likely to correspond to the naive stage of embryogenesis, which is characterized by the expression of naive pluripotency markers and the ability to form ES cells.  The study also investigated the adhesive properties of finch and chick blastodermal cells in culture and found that finch cells were more sensitive to the substrate type, while chick cells were more robust. The researchers used a variety of substrates, including gelatin, collagen, and E-cadherin, and found that finch cells grew optimally on human placenta type IV collagen-coated surfaces.  
Overall, the study provides new insights into the early development of the zebra finch and its potential use as a model organism for studying embryogenesis and stem cell biology. The findings suggest that the finch blastoderm at oviposition is a valuable resource for studying the naive stage of embryogenesis and the derivation of ES cells. \\ \hline
50 & The study aimed to investigate the early development of the zebra finch embryo and its potential use for comparative stem cell studies and genome engineering. The researchers used a combination of morphological and molecular characterization to study the finch embryo at oviposition, which is equivalent to the blastocyst stage in mouse embryos. They found that the finch embryo at oviposition is similar to the mouse embryo at E4-E4.5, a stage that has been termed the naive state of embryogenesis.  The researchers used quantitative PCR (Q-PCR) to compare the expression levels of various genes associated with early embryonic development and pluripotency in the finch and chick blastoderms at oviposition. \texttt{[...]}
% They found that the finch blastoderm at oviposition showed higher levels of expression of naive pluripotency markers, such as Fbxo15, Prdm14, and Tbx3, compared to the chick blastoderm. They also found that the finch blastoderm at oviposition showed higher levels of expression of genes associated with the naive state of pluripotency, such as Nrob1 and Esrrb. 
% The researchers then cultured the finch and chick blastoderms at oviposition in the presence of LIF and the MEK inhibitor, PD0325901, and found that the finch blastodermal cells formed aggregates that showed strong AP activity and expressed Nanog and PouV. In contrast, the chick blastodermal cells did not show AP activity or express Nanog and PouV, even after extended staining incubation times.  The researchers concluded that the finch blastoderm at oviposition is likely to correspond to the naive stage of embryogenesis, which is a stage that has been conserved amongst amniotes. They also suggested that the ability to form epithelial-type adhesive contacts is present but inhibited at early/mid-blastula stages, and that the expression of E-cadherin represents a metastable, transitional state during mesenchymal-epithelial transition that is permissive for ES cell derivation.  
Overall, the study provides new insights into the early development of the zebra finch embryo and its potential use for comparative stem cell studies and genome engineering. The findings suggest that the finch blastoderm at oviposition is a valuable model system for studying the naive state of embryogenesis and the mechanisms of ES cell derivation. \\ \hline
75 & The study aimed to investigate the early development of the zebra finch embryo and its potential use for comparative stem cell studies and genome engineering. The researchers used a combination of morphological and molecular characterization to study the finch embryo at oviposition, which is equivalent to the blastocyst stage mouse embryo. They found that the finch blastoderm at oviposition was similar to the mouse embryo at the same stage, with a multi-layered structure and the presence of epiblast and hypoblast precursors.  The researchers also found that the finch blastoderm at oviposition expressed markers of pluripotency, including Nanog, PouV, and Dnmt3b, which are associated with the naive state of pluripotency. \texttt{[...]}
% They also found that the finch blastoderm at oviposition showed higher levels of expression of naive pluripotency markers, such as Fbxo15, Prdm14, and Tbx3, compared to the chick blastoderm at oviposition.  
% The researchers then cultured finch and chick blastodermal cells in the presence of LIF and the MEK inhibitor, PD0325901, and found that finch blastodermal cells formed aggregates that showed strong AP activity and expressed Nanog and PouV, while chick blastodermal cells did not. They also found that finch blastodermal cells were insensitive to MEK inhibition, while chick blastodermal cells required MEK for the expression of pluripotent markers.  The study suggests that the finch blastoderm at oviposition is likely to correspond to the naive stage of embryogenesis, which is a stage of development that is characterized by the presence of a multi-layered structure and the expression of pluripotency markers. The study also suggests that the finch embryo may be a useful model for studying the early development of amniotes and for understanding the mechanisms of pluripotency and stem cell maintenance.  
Overall, the study provides new insights into the early development of the zebra finch embryo and its potential use for comparative stem cell studies and genome engineering. The study also highlights the importance of the finch embryo as a model system for studying the mechanisms of pluripotency and stem cell maintenance.  \texttt{[...]}
% Readability: 75/100  Note: The readability score is based on the Flesch-Kincaid Grade Level test, which measures the complexity of written text. A score of 75/100 indicates that the text is written at a level that is accessible to readers with a high school education, but may require some effort to understand for readers with a lower level of education. 
\\ \hline
100 & The study aimed to investigate the early development of the zebra finch (Taeniopygia guttata) and its potential use as a model organism for studying embryogenesis and stem cell biology. The researchers used a combination of morphological and molecular characterization techniques to study the finch embryo at oviposition, including quantitative RT-PCR, immunohistochemical staining, and in situ hybridization.  \texttt{[...]}
% The results showed that finch embryos are laid at a stage equivalent to the blastocyst stage in mouse embryos, with a multi-layered blastoderm that has not yet epithelialized. The epiblast and hypoblast precursors were already molecularly specified at this stage, and the primordial germ cell fate was marked by immunoreactivity to SSEA-1 and EMA-1 antibodies. 
 % The researchers also found that the finch blastoderm at oviposition expressed markers of pluripotency, including Nanog, PouV, and Dnmt3b, which are associated with the naive state of pluripotency. In contrast, the chick blastoderm at oviposition showed lower levels of expression of these markers.  The study also investigated the effect of MEK inhibition on the expression of pluripotency markers in finch and chick blastodermal cells. The results showed that finch blastodermal cells cultured with MEK inhibitor retained markers of pluripotency, including AP activity and expression of Nanog and PouV, whereas chick blastodermal cells did not.  The researchers concluded that the finch embryo at oviposition is likely to correspond to the naive stage of embryogenesis, which is characterized by the expression of pluripotency markers and the absence of epithelialization. 
They also suggested that the finch embryo may be a useful model organism for studying embryogenesis and stem cell biology, particularly in the context of avian development.  Overall, the study provides new insights into the early development of the zebra finch and its potential use as a model organism for studying embryogenesis and stem cell biology. The findings have implications for our understanding of the evolution of germ layer formation in amniotes and the development of new technologies for manipulating avian development. \\ \hline
\end{tabular}
\caption[Example Outputs from In-Context Learning at each requested Readability Level.]{Example Outputs from In-Context Learning at each requested Readability Level (RL). Readability level is on a scale of $0$ to $100$, with $100$ being the most readable.}\label{tab:icl-ex-outputs}
\end{table*}

\begin{table*}[]
\centering
\scriptsize
\begin{tabular}{l|p{.89\textwidth}}
\textbf{RL} & \textbf{CV Output} \\ \hline
-1.0 & The wild fish samples were collected in 2015 in the Gonarezhou National Population in Zimbabwe and Mozambique. Intestines were collected at each location and preserved in pure ethanol. Sampling locations coordinates are listed in Figure 2—source data 1. Fish (GRZ strain) used for microbiota analysis and scored for survival were individually housed from week 4 post-hatching in single 2.8L tanks connected to a water recirculation system receiving 12 hr of light and 12 hr of dark every day. Water temperature was set to 28°C and fish were fed blood worm larvae and brine shrimp nauplii twice a day during the week and once a life span in the wild fish populations, we identified a core microbiota that is conserved between wild and captive TK, with a significant reduction in gut bacterial richness in old fish. \\ \hline
-0.5 & The turquoise killifish (TK) is a naturally short-lived vertebrate model system, characterized by a broad spectrum of aging phenotypes, including cancer, neurodegeneration, and behavioral decline. We show that TK have a complex gut microbial taxonomic diversity similar to other vertebrate aging model organisms, including zebrafish, mice, and humans. The core microbiota is conserved between wild and captive TK, with a significant reduction in gut bacterial richness during aging. Acute gut microbiota transfer in the context of normal aging significantly prolongs life span in a vertebrate, becoming a novel candidate life span enhancing intervention. Our results indicate that improving the ecological diversity of the GM in old individuals helps to restore health and prolongs life span. Our approach could provide a key to slowing aging and retarding the onset of age-associated diseases by specifically targeting the GM. The wild fish samples were collected in 2015 during an expedition in the Gonarezhou National Park in Zimbabwe and Mozambique. Intestines were collected at each location and preserved in pure ethanol. Sampling locations coordinates are listed in Figure 2—source data 1. \\ \hline
0.0 & The study investigates the relationship between gut microbiota and aging in the turquoise killifish, a short-lived vertebrate model organism. The researchers found that the gut microbiota of young fish is more diverse and complex than that of old fish, with a higher abundance of beneficial bacteria such as Firmicutes and Bacteroidetes. In contrast, old fish have a less diverse gut microbiota with a higher abundance of pathogenic bacteria such as Proteobacteria. The researchers also found that the gut microbiota of old fish is associated with inflammation and immune responses, while the gut microbiota of young fish is associated with healthy gut function and cell proliferation. \texttt{[...]}
% The study suggests that the gut microbiota plays a key role in aging and that manipulating the gut microbiota may be a potential strategy for promoting healthy aging. The researchers also found that transferring young gut microbiota to middle-aged fish can prolong their lifespan and improve their health, suggesting that the gut microbiota may be a key factor in determining lifespan. Overall, the study provides new insights into the relationship between gut microbiota and aging and highlights the potential of manipulating the gut microbiota as a strategy for promoting healthy aging. 
The study used a combination of 16S rRNA gene amplicon sequencing and RNA sequencing to analyze the gut microbiota and transcriptome of young and old turquoise killifish. The researchers found that the gut microbiota of young fish is more diverse and complex than that of old fish, with a higher abundance of beneficial bacteria such as Firmicutes and Bacteroidetes. In contrast, old fish have a less diverse gut microbiota with a higher abundance of pathogenic bacteria such as Proteobacteria. \\ \hline
0.5 & The study looked at the relationship between the gut microbiome and aging in the turquoise killifish. The researchers found that as the fish get older, their gut microbiome changes and becomes less diverse. They also found that if they took the gut microbiome from a younger fish and put it into an older fish, it could help to extend the older fish's lifespan. The researchers also looked at the genes that are expressed in the gut of older fish and found that they are different from the genes that are expressed in the gut of younger fish. \texttt{[...]}
% They also found that if they took the gut microbiome from a younger fish and put it into an older fish, it could help to change the way that the older fish's genes are expressed. The study suggests that the gut microbiome could be an important factor in how long we live, and that it could be possible to use it to help to extend our lifespan. The researchers also found that the gut microbiome of the turquoise killifish is similar to the gut microbiome of humans, and that it could be possible to use the turquoise killifish as a model to study how the gut microbiome affects our health and how it could be used to help to extend our lifespan. The study was conducted by a team of researchers from the University of Bologna and the University of California, San Francisco. The researchers used a technique called 16S rRNA gene sequencing to study the gut microbiome of the turquoise killifish. 
They also used a technique called RNA sequencing to study the genes that are expressed in the gut of the fish. The researchers also looked at the lifespan of the fish and how it was affected by the gut microbiome.\\ \hline
1.0 & The researchers are looking at how to live longer by looking at what is in the gut.  They are looking at the gut of a fish that lives for about 10 years.  They want to see if they can help the fish live longer by changing what is in the gut.  They are also looking at what is in the gut of people.  They want to see if they can help people live longer by changing what is in their gut.  They are also looking at what is in the gut of mice.   They want to see if they can help mice live longer by changing what is in their gut. \\ \hline
\end{tabular}
\caption[Example outputs from Control Vector method at each requested Readability Level.]{Example outputs from Control Vector method at each requested Readability Level (RL). Readability level is on a scale of $-1.0$ to $1.0$, with $1.0$ being the most readable.}\label{tab:cv-ex-outputs}
\end{table*}

